\section{Which root contrasts require an interaction in the three training grammars?}
\label{sec:rule-audit}

The parity example has a root contrast whose conditional mean given either
child alone is zero. We can check that property exactly in the three
quenched grammars used in the root-classification experiments. For each of
the \(n=16\) root symbols, enumerate its \(m=4\) legal top productions. A
row \(e\) records a root label \(j_e\) and an ordered child pair
\((u_e,v_e)\). Under the uniform root and uniform production law, the
\(nm=64\) rows have equal probability. The calculation uses the actual
hidden child symbols and asks what information is available even with
perfect access to them.

Let \(E\) be the \(64\times16\) one-hot label matrix and let \(D\) be the
\(64\times32\) matrix whose row has a one in left column \(u_e\) and
another in right column \(16+v_e\), using zero-based indices. A root contrast
\(c\in\mathbb R^{16}\) assigns the value \(y_e=(Ec)_e=c_{j_e}\) to
each row. Every prediction of the form \(f(u_e)+g(v_e)\) lies in
\(\operatorname{col}D\), and every vector in that column space has such a
representation. Centered root contrasts satisfy \(\sum_j c_j=0\).

\paragraph{Two dimension counts.}
The vector \(D^\top Ec\) contains the sums of \(y_e\) within each left
symbol and within each right symbol. On a symbol of positive probability,
its sum vanishes exactly when its conditional mean vanishes. Let us call
a contrast \emph{purely interactive} when both sets of marginal signals
vanish. The kernel of \(D^\top E\) consists of exactly those contrasts.
They are already centered because constants belong to \(\operatorname{col}D\).
At the opposite extreme \(Ec\) itself may be additive on the legal
productions. Since \(E^\top E=m I\), the map from contrasts to edge
values is injective. The dimension formula for the intersection of two
subspaces gives
\begin{align}
 d_{\rm pure} &= n-\operatorname{rank}(D^\top E),\\
 d_{\rm additive,centered}
 &= n+\operatorname{rank}D-\operatorname{rank}[D\ E]-1 .
\end{align}
The last subtraction removes the shared constant. Both ranks are computed
over the rationals from integer matrices.

\paragraph{How much can an additive predictor explain?}
Let \(P_D=DD^\dagger\), the orthogonal projection onto additive functions
on these rows. For a nonzero centered contrast, the fraction of its
population variance explained by its best additive predictor is
\[
 \eta(c)=\frac{\|P_D Ec\|^2}{\|Ec\|^2}.
\]
Indeed, the uniform row weights cancel in the ratio, and the projection
minimizes the squared residual. If \(H\) has orthonormal columns spanning
\(\mathbf 1^\perp\) in label space, then \(Q=EH/\sqrt m\) has orthonormal
columns in row space. The minimum and maximum of \(\eta\) are the extreme
eigenvalues of \(Q^\top P_DQ=(P_DQ)^\top(P_DQ)\). These fractions allow
both children's separate contributions to enter in the same additive fit,
including any correlations between the two child symbols.

\begin{center}
\begin{tabular}{c|ccc|cc}
Seed & \(\operatorname{rank}D\) & \(\operatorname{rank}(D^\top E)\)
& \(\operatorname{rank}[D\ E]\) & \(\min\eta\) & \(\max\eta\)\\ \hline
0 & 30 & 16 & 45 & 0.063505 & 0.927720\\
1 & 31 & 16 & 46 & 0.108409 & 0.821812\\
2 & 31 & 16 & 46 & 0.063091 & 0.909308
\end{tabular}
\end{center}

For each seed both dimension counts are zero. Every nonzero centered root
contrast has some marginal signal, so a pure parity contrast cannot be
selected from the sixteen label values in these tables. Every such contrast
also retains a nonzero residual after its best additive fit. For seed 0,
the least additive direction retains \(93.65\%\) of its variance in that
residual. That makes it a candidate for studying how learning the children
changes the root's access to information about their joint state.
Residualizing \(Ec\) against \(D\) would produce an exactly pure edge
target, but its values can differ across productions of the same root.
That intervention would change the prediction task.

To reproduce the calculation from the repository root, run
\begin{quote}
\small\ttfamily python3 -m studies.rhm.child\_feature\_enabling.grammar\_audit
\end{quote}
It enumerates the root tables from \texttt{studies.rhm.common.GRAMMAR}
with seeds 0, 1, and 2. The accompanying JSON records the tables, all
fifteen eigenvalues per seed, the exact ranks, the grammar settings, and
hashes of the grammar source files. The code checks the projection spectrum
against an independent principal-angle calculation. A planted binary parity
table and a table in which the root equals the left child check the pure
and additive endpoints, respectively.
